% ! TeX program = lualatex
% A function and its inverse as reflections in the line y = x.  Used by the
% midterm review slides.
%
%   page 1 -- y = f(x), the line y = x, and the point (2, 3) on the curve
%   page 2 -- the same picture with f^{-1} and the reflected point (3, 2)
%
% The curve is f(x) = x^3/4 + 1, which is increasing (so it has an inverse),
% is nowhere near symmetric in y = x (so the reflection is visibly a second
% curve), and passes through (2, 3), (0, 1) and (-2, -1), so the swap can be
% read off the grid in whole numbers.
%
% The inverse is drawn parametrically as (f(t), t) rather than from a formula
% for f^{-1}.  That is the same exchange of x and y the slide is about, and it
% sidesteps plotting a cube root through negative numbers.
\documentclass[multi=tikzpicture,crop]{standalone}

\input{../typography.tex.preamble}
\input{../tikz.tex.preamble}

\pgfplotsset{inverse reflection/.style={
    basic curve,
    enlargelimits=false,
    % extra room on the right for the y = f^{-1}(x) label
    xmin=-2.9, xmax=5.3,
    ymin=-2.9, ymax=4.4,
    width={210pt},
  }}

% y = x, the mirror, plus the original curve and the point we follow.
\newcommand{\originalcurve}{%
  \addplot[dashed] coordinates {(-2.9,-2.9) (4.4,4.4)};

  \addplot[main, very thick, domain=-2.2:2.2] {(x^3)/4 + 1};

  \draw[main, fill=main] (axis cs:2,3) circle (0.07);
}

\begin{document}

% ---------------------------------------------------------------------
% Page 1 -- the function and the mirror.
\begin{tikzpicture}
  \begin{axis}[inverse reflection]
    \originalcurve{}
  \end{axis}
\end{tikzpicture}

% ---------------------------------------------------------------------
% Page 2 -- the reflection, drawn as (f(t), t).
\begin{tikzpicture}
  \begin{axis}[inverse reflection]
    \originalcurve{}

    \addplot[supp, very thick, variable=\t, domain=-2.2:2.2, samples=100]
    ({(t^3)/4 + 1}, {t});
    % the one curve under both of its names, each sitting above it
    \node[main, fill=white, inner sep=1pt] at (axis cs:2,4)
    {\(x = f(y)\)};
    \node[supp, fill=white, inner sep=1pt] at (axis cs:4,2.5)
    {\(y = f^{-1}(x)\)};

    \draw[supp, fill=supp] (axis cs:3,2) circle (0.07);
  \end{axis}
\end{tikzpicture}
\end{document}
